\documentclass[11pt,a4paper]{article}
\usepackage[margin=27mm]{geometry}
\usepackage{amsmath,amssymb,amsthm,hyperref}
\newtheorem{theorem}{Theorem}
\newtheorem{lemma}[theorem]{Lemma}
\newtheorem{corollary}[theorem]{Corollary}
\newcommand{\gen}[1]{\langle #1\rangle}
\newcommand{\Z}{\mathbb Z}
\title{Finite quotients preserving spherical residues}
\author{Human Compression project}
\date{}
\begin{document}\maketitle

\begin{abstract}
For a Coxeter system with finite maximal parabolics, a finite set of
products of three parabolics detects every failure of flag lifting in
a quotient. Avoiding this set preserves the entire incidence structure
of each proper residue, as well as the intersection property.
Reduction of the faithful geometric representation modulo a suitable
prime produces such quotients of arbitrarily large order whenever the
Coxeter group is infinite. This gives a common proof of the finite
existence assertions in Section~6 of the fixed manuscript, including
the diagrams which are not strings. It does not determine the smallest
quotient or verify the groups in the source's finite table.
\end{abstract}

\section{A finite obstruction set}

Let $(W,S)$ be a Coxeter system of rank $r\ge3$, with
$S=\{s_i:i\in I\}$. Put $W_J=\gen{s_j:j\in J}$ for $J\subseteq I$,
and write $P_i=W_{I\setminus\{i\}}$. Assume that every $P_i$ is finite.
We use the standard Coxeter complex: its objects of type $i$ are the
left cosets $wP_i$, and two objects of different types are incident
when their cosets intersect. We retain the classical Coxeter facts
that $W_J\cap W_K=W_{J\cap K}$, that this incidence system is a thin,
residually connected geometry, and that $W$ acts regularly on its
chambers. In particular every pairwise incident set extends to a
chamber. These are inputs about the universal complex, not assumptions
about a proposed quotient.

Define the finite set
\[
 \mathcal F=\biggl(\bigcup_{i,j,k\in I} P_iP_jP_k\biggr)
                 \setminus\{1\}.
\]
Repetitions among $i,j,k$ are allowed. Since every subgroup contains
the identity, this set also contains all nonidentity elements in any
$P_i$ or $P_iP_j$.

\begin{theorem}\label{lifting}
Let $N\triangleleft W$ satisfy $N\cap\mathcal F=\varnothing$, and
let $q:W\to Q=W/N$. The marked group
$(Q,(q(s_i))_{i\in I})$ has the full intersection property.
Its coset incidence system with maximal subgroups $q(P_i)$ is a
regular hypertope. The residue of its base object of type $i$ is
isomorphic, as a typed incidence geometry, to the residue of $P_i$
in the universal Coxeter complex. Thus every proper residue remains
universal and spherical. The quotient map preserves every bond label.
\end{theorem}

\begin{proof}
First, $q$ is injective on every $P_i$. More generally, let $J,K$
be proper subsets of $I$ and choose $i\notin J$, $j\notin K$.
If $a\in W_J$, $b\in W_K$ and $q(a)=q(b)$, then
$ab^{-1}\in N\cap P_iP_j$, so $a=b$. The intersection property
in $W$ therefore gives
\[
 q(W_J)\cap q(W_K)=q(W_{J\cap K}).
\]
When either subset is all of $I$, this equality is immediate.
This proves the full intersection property, including the trivial
intersection of all maximal marked parabolics. Every subgroup
generated by two marks is contained in a maximal parabolic, since
$r\ge3$. Its injection proves that no mark disappears and that all
pairwise product orders are unchanged.

We next prove the required flag lifting, rather than inferring it
from these subgroup equalities. Let a nonempty flag in the quotient
contain an object of type $i$. Translate it so that this object is
$q(P_i)$. Every other object, of type $j$, is incident with
$q(P_i)$, and hence has a representative
\[
                         q(a_j)q(P_j),\qquad a_j\in P_i.
\]
Choose such a representative separately for every type in the flag.
The corresponding objects $a_jP_j$ are all incident with $P_i$.
For distinct types $j,k$ in the flag, quotient incidence says
\[
 q(a_j^{-1}a_k)\in q(P_jP_k).
\]
Choose $b\in P_j$, $c\in P_k$ with
$q(a_j^{-1}a_k)=q(bc)$. Then
\[
        a_j^{-1}a_kc^{-1}b^{-1}\in N\cap P_iP_kP_j.
\]
The avoidance hypothesis makes this element the identity. Consequently
$a_j^{-1}a_k=bc\in P_jP_k$, which is exactly the incidence of
$a_jP_j$ and $a_kP_k$. Thus the chosen objects form an actual flag
upstairs. They extend to a chamber of the Coxeter complex. Its image
is a chamber containing the prescribed quotient flag.
The same argument proves that every quotient flag is a translate of
a subflag of the base chamber: lift, extend and use the chamber action
of $W$. In particular $Q$ is flag transitive.

Here is the asserted residue isomorphism in detail. An object of type
$j\ne i$ incident with $P_i$ has the form $aP_j$ for $a\in P_i$.
Mapping it to $q(a)q(P_j)$ is surjective onto the corresponding
quotient residue. If $a,b\in P_i$ give the same image object, then
$q(a^{-1}b)\in q(P_j)$. The already proved intersection equality
and injectivity on $P_i$ imply
$a^{-1}b\in P_i\cap P_j$, so the objects upstairs coincide.
This proves injectivity on objects. The preceding triple-product
calculation proves that incidence of any two image objects reflects
incidence upstairs. Incidence is preserved automatically. Hence the
map is an isomorphism of the complete residue geometries, not merely
of their groups or their chamber sets.

Every nonempty proper residue in the quotient is therefore a residue
of a finite spherical Coxeter complex. It is thin and, when its rank
is at least two, connected. The whole quotient geometry is connected
as well: its chambers are the translates of the base chamber, with
trivial chamber stabilizer by the intersection property. Multiplying
successively by the marks gives galleries joining these chambers,
because the marks generate $Q$. A gallery gives connectedness of
the incidence graph in rank at least two. Thinness and residual
connectedness now hold in all required ranks. Finally, in this geometry, a type-preserving automorphism fixing a
chamber fixes its adjacent chambers and then every chamber, by the
gallery connectedness just proved. Thus the full
type-preserving automorphism group acts freely on chambers. Since
$Q$ already acts transitively and freely, it is that full group.
\end{proof}

The third factor has a specific role: two objects represented inside
one fixed parabolic may become incident after quotienting. Pairwise
parabolic injection alone does not rule this out. The argument does
not claim that the displayed avoidance condition is necessary, or
that it is economical for finding a small quotient.

\section{Separating finitely many matrices}

We give the elementary finite-quotient argument needed here. It also
specifies what is retained from Coxeter theory.

\begin{lemma}\label{matrices}
Let $K$ be a number field with ring of integers $\mathcal O_K$,
and let $G\le\operatorname{GL}_d(\mathcal O_K)$. For every finite
set $A\subset G\setminus\{1\}$ there is a homomorphism from $G$
to a finite group whose kernel is disjoint from $A$.
\end{lemma}

\begin{proof}
For every $g\in A$, choose a nonzero entry $a_g$ of $g-I$.
Its algebraic norm $\operatorname{Norm}_{K/\mathbb Q}(a_g)$ is a
nonzero integer: it is the determinant of multiplication by $a_g$
on any integral basis of $\mathcal O_K$. Choose a rational prime
$p$ not dividing the product of these finitely many integers; for
empty $A$ any prime will do. Modulo $p$, each multiplication
matrix remains invertible. In particular $a_g$ has nonzero image
in the finite ring $\mathcal O_K/p\mathcal O_K$.

Reduction of matrices is a homomorphism
\[
 G\longrightarrow
 \operatorname{GL}_d(\mathcal O_K/p\mathcal O_K).
\]
Inverses remain inverses after reduction. The target is finite,
since the additive group of the coefficient ring has
$p^{[K:\mathbb Q]}$ elements. The chosen entry of every $g-I$
survives, so none of the elements of $A$ maps to the identity.
\end{proof}

For a finite-rank Coxeter system the standard geometric representation
has, in its simple-root basis, the matrices
\[
 s_i(e_i)=-e_i,\qquad
 s_i(e_j)=e_j+2\cos(\pi/m_{ij})e_i\quad(j\ne i).
\]
The usual convention for $m_{ij}=\infty$ replaces the coefficient
by $2$. We retain the faithfulness theorem for this representation.
Each finite-label coefficient is an algebraic integer:
$2\cos(\pi/m)=\zeta_{2m}+\zeta_{2m}^{-1}$.
There are only finitely many labels, so all coefficients lie in the
ring of integers of a single cyclotomic number field. The generators
are involutions and have determinant $-1$. Thus both the matrices
and their inverses have integral coefficients, and Lemma~\ref{matrices}
applies. This proves the particular residual-finiteness assertion used
below without invoking a normal-subgroup classification.

\begin{corollary}\label{existence}
Let $(W,S)$ be an infinite Coxeter system of finite rank at least
three, with every maximal standard parabolic finite. For every
positive integer $M$ there is a finite regular hypertope with this
diagram, with universal spherical proper residues, and with at least
$M$ chambers.
\end{corollary}

\begin{proof}
Choose $M$ distinct elements $w_1,\ldots,w_M$ of $W$ and enlarge
the finite set $\mathcal F$ by the nonidentity elements
$w_a^{-1}w_b$ for $a\ne b$. Apply the faithful representation
and Lemma~\ref{matrices} to separate this entire set simultaneously.
The image group $Q$ is finite, and its kernel satisfies
Theorem~\ref{lifting}. The images of the $w_a$ are distinct,
so $|Q|\ge M$. The regular chamber action identifies $|Q|$ with
the number of chambers of the resulting hypertope.
\end{proof}

This proof is effective in an elementary sense. The finite parabolics
can be enumerated in the exact representation, their triple products
formed, and a nonzero matrix entry selected for every forbidden
element. The norm calculation then supplies a prime to avoid.
The finite matrix image can itself be enumerated. These steps need
not be computationally practical; no small-order bound is asserted.
Taking unbounded $M$ gives infinitely many pairwise nonisomorphic
finite hypertopes, rather than just many maps onto the same object.

\section{Source consequences and attribution}

Section~6 of the fixed manuscript~\cite{source} recalls infinitely
many finite maps and hypermaps for the compact hyperbolic triangle
diagrams. It also recalls, from the classical polytope literature,
infinitely many finite locally spherical polytopes for the string
diagrams in ranks four and five. Corollary~\ref{existence} supplies
one proof for these assertions. Indeed a compact hyperbolic Coxeter
simplex has finite proper parabolics, and its reflection group is
infinite. The same argument applies to the nonstring compact
hyperbolic diagrams in the source's Table~2. We retain that diagram
classification; the corollary does not provide a new proof that the
list is exhaustive. For a string diagram the acquired string
coset correspondence identifies the resulting hypertope with a
regular abstract polytope.

The conclusion uses actual universal spherical residues. The source
also uses ``spherical'' for the weaker finite-diagram condition,
which may permit projective residues. Our existence result satisfies
the stronger condition. It cannot be used to rule out smaller
examples with projective residues. Nor does it identify a quotient
group with any particular extension listed in Table~3, prove a
minimum there, or verify a quoted exceptional construction. Those
remain separate finite or attributed inputs.

The finite existence results for string diagrams are prior work,
as are residual finiteness of Coxeter groups and the standard
Coxeter complex. The contribution of this component is the explicit
triple-product argument linking a finite separating set to full flag
lifting and residue incidence, followed by the common application.
No claim of historical priority is made for this argument. The
earlier project note \emph{Finite quotients with prescribed parabolic
intersections} treats different hypotheses involving separability
and endomorphisms; none of its unproved filling assertions is used
here. This component makes no claim to resolve a source conjecture,
complete the entire L033 manuscript, provide a Lean formalization,
or establish a numerical compression factor.

\begin{thebibliography}{9}
\bibitem{source} M.~E.~Fernandes, D.~Leemans and A.~Ivi\'c Weiss,
\emph{An exploration of locally spherical regular hypertopes},
fixed sixteen-page revised manuscript, catalogue L033,
Section~6 and Table~2. Its classical existence citations and the
finite table are distinguished above.
\bibitem{acquired} Human Compression project,
\emph{Residue stabilizers and distinguished generators};
\emph{Kernels and intersections in quotients of hypertopes}.
The present proof retains the universal Coxeter complex and its
standard intersection property, and proves the quotient interfaces
explicitly.
\end{thebibliography}
\end{document}
